\documentclass[10pt,a4paper]{amsart}
\usepackage[margin=19mm]{geometry}
\usepackage{amsmath,amssymb,mathtools}
\usepackage[scr=rsfs,cal=euler]{mathalfa}
\usepackage{xfrac}
\usepackage[unicode,hidelinks,bookmarks=false]{hyperref}
\newcommand{\ie}{i.e.\ }
\newcommand{\cf}{cf.\ }
\newcommand{\resp}{respectively\ }
\newcommand{\Subsection}{\subsection}
\providecommand{\nobreakdash}{\nobreak\hbox{-}}
\setlength{\emergencystretch}{2em}
\multlinegap0pt
\usepackage{amssymb}
\usepackage[shortlabels]{enumitem}
\usepackage{tikz}
\newtheorem{theorem}{Theorem}
\newtheorem{lemma}{Lemma}
\title{A note on words having the same image on finite groups}
\author{Shrinit Singh}
\date{Source: arXiv:2407.00789v2 (CC BY 4.0)}
\begin{document}
\maketitle

\noindent\emph{Source and licence.} A note on words having the same image on finite groups. Version arXiv:2407.00789v2 (CC BY 4.0), distributed by arXiv under the Creative Commons Attribution 4.0 International licence, http://creativecommons.org/licenses/by/4.0/, as recorded in the arXiv licence metadata for this exact version and shown on the abstract page https://arxiv.org/abs/2407.00789v2. Full licence notice: \texttt{licenses/arxiv-2407.00789-CC-BY-4.0.txt}.\par\smallskip

\noindent\textit{Editorial scope.} Counted unit: the article's theorem main (source0.tex lines 213--221). The article's complete original proof of it is delivered with it (source0.tex lines 252--271, one \texttt{proof} environment of 20 lines, which the article places away from the statement). No mathematics is written, repaired or re-derived: the statement and the proof are the article's own lines, in the article's own notation, and no retained line is substituted or re-flowed.  Retained uncounted as the source-local context the proof needs: lines 226--228.  Nothing was omitted from any retained span, so the package carries no omission record.  No retained line contains a citation, so the wrapper carries no bibliography and no bibliography entry.  No retained line contains a figure environment, an image-inclusion command or drawing code, so no image is delivered and the package declares no asset dependency.  Editorial formatting only, in addition: an AMS \texttt{amsart} preamble that restates the article's own declarations and macros used by the retained lines, namely the environments \texttt{theorem}, \texttt{lemma} on separate counters, together with the standard packages those lines need.  The retained environments print in this document as Theorem 1, Lemma 1 in order of appearance, whereas the article prints the same items with the numbering of its own full document.  The complete native source of the article as distributed in the arXiv e-print for this exact version is bundled unchanged as \texttt{sources/161590/source0.tex}.
\begin{lemma}\label{promon}
    Let $P$ be a finitely generated group. Then $\mathrm{End} (\widehat{P}) = \varprojlim_{J} \mathrm{End}(P/J)$, where the inverse limit is taken over all $J \leq P$ such that $J = J_P(G)$ for some finite group $G.$
\end{lemma}

\begin{theorem}\label{main}
    Let $P$ be a finitely generated group and $\gamma_1, \gamma_2 \in P.$ Then the following are equivalent:
    \begin{enumerate}
        \item $i(\gamma_1)$ and $i(\gamma_2)$ are endomorphic equivalent  in $\widehat{P}$.
        \item For every finite group $G$, and every $g \in G$, the set $\mathrm{Hom}_{\gamma_1,g}(P,G)$ is non-empty if and only if $\mathrm{Hom}_{\gamma_2,g}( P,G)$ is non-empty.
        \item $\gamma_1J$ and $\gamma_2J$ are endomorphic equivalent in $P/J,$ where $J= J_P(G)$ for every finite group $G.$
        \item For every $N \unlhd _\mathrm{f.i.}P,$ there exists $J \unlhd_\mathrm{f.i.} P$ with $J \leq N$ such that $\gamma_1J$ and $ \gamma_2J$ are endomorphic equivalent in $P/J.$
    \end{enumerate}
\end{theorem}

\begin{proof}[Proof of Theorem \ref{main}]
\mbox{}
\begin{description}
    \item[$1 \implies 2$] The proof follows from the fact that $\mathrm{Hom}(P,G)$ and $\mathrm{Hom}(\widehat{P},G)$, for every finite group $G$ are in one-one correspondence and we have $| \mathrm{Hom}_{\gamma,g} (P,G)| = |\mathrm{Hom}_{i(\gamma),g}(\widehat{P}, G)|.$ Hence, $$ \vert \mathrm{Hom}_{\gamma_1,g} (P,G)| = |\mathrm{Hom}_{i(\gamma_1),g}(\widehat{P}, G)|  = |\mathrm{Hom}_{i(\gamma_2),g}(\widehat{P}, G)| = |\mathrm{Hom}_{\gamma_2,g} (P,G)|.$$
    \item [$2 \implies 3$] Let $J = J_P(G)$ for some finite group $G$. As $\mathrm{Hom}_{\gamma_1, \gamma_1J}(P, P/J)\neq \emptyset \implies \mathrm{Hom}_{\gamma_2, \gamma_1J}(P,P/J) \neq \emptyset$. Let $f \in \mathrm{Hom}_{\gamma_2, \gamma_1J}(P,P/J).$ Then $J \subseteq \ker(f)$ because of the choice of $J.$ Hence, there exists an induced homomorphism $\Tilde{f}: P/J \longrightarrow P/J$ such that $(\tilde f) (\gamma_2J) = \gamma_1J.$ Similarly, we can show the other way. Hence, $\gamma_1J$ and $\gamma_2J$ are endomorphic equivalent in $P/J$.
    \item [$3 \implies 4$] For every $N \unlhd_{\mathrm{f.i.}}P,$ we can take $J = J_P(P/N).$ 
    \item [$4 \implies 3$] Let $J = J_P(G)$ for some finite group $G.$ By assumption, there exists a subgroup $J' \le J$ such that $J' \unlhd_{\mathrm{f.i.}} P$ with $\gamma_1J'$ and $\gamma_2J'$ endomorphic equivalent. But the image of $J$ in $P/J'$ is $J_{P/J'}(G),$ so this image is fully invariant, hence, every endomorphism of $P/J'$ induces an endomorphism in $P/J$. Hence, we conclude that $\gamma_1J, \gamma_2J$ are endomorphic equivalent in $P/J.$
    \item [$3 \implies 1$] For each $J = J_P(G)$, define
$$
S_J = \left\{ \alpha \in \mathrm{End}(P/J) \mid \alpha(\gamma_1 J) = \gamma_2 J \right\}.
$$
Each $S_J$ is non-empty by hypothesis and closed in the finite monoid $\mathrm{End}(P/J)$ (endowed with the discrete topology). The canonical projections $$\pi_J: \varprojlim_{J} \mathrm{End}(P/J) \to \mathrm{End}(P/J)$$ are continuous. Thus, $\pi_J^{-1}(S_J)$ is closed in $\varprojlim_{J} \mathrm{End}(P/J)$ for each $J$.

For any finite collection $\{J_1, \dots, J_n\}$, directedness yields $J \subseteq \bigcap_{i=1}^n J_i$ with $J = J_P(G)$. For $\alpha \in S_J$, full invariance of $J_i/J$ in $P/J$ ensures $\alpha$ induces $\alpha_{J_i} \in S_{J_i}$ which is also the image of $\alpha$ in $\mathrm{End}(P/J_i)$. In other words, an element of $\varprojlim_{J} \mathrm{End}(P/J)$ defined by $\alpha$ at $J^{\text{th}}$ coordinate forces the $J_i^{\text{th}}$ coordinate to be $\alpha_{J_i}$ for each $1 \leq i \leq n$. Hence, the compatibility condition implies that $\pi_{J}^{-1}(\alpha) \subseteq \pi_{J_i}^{-1}(\alpha_{J_i})$ for each $1 \leq i \leq n$. Thus, $\bigcap_{i=1}^n \pi_{J_i}^{-1}(S_{J_i}) \neq \emptyset$. So the family $\{\pi_J^{-1}(S_J)\}_J$ has the finite‐intersection property. Compactness of $\varprojlim_{J} \mathrm{End}(P/J)$ then gives:
$$
\bigcap_J \pi_J^{-1}(S_J) \neq \emptyset.
$$
Take $(\alpha_J)_J$ in this intersection. Then $\alpha_J(\gamma_1 J) = \gamma_2 J$ for all $J$. By Lemma \ref{promon}, there exists $\widehat{\alpha} \in \mathrm{End}(\widehat{P})$ mapping $i(\gamma_1) = (\gamma_1 J)_J$ to $i(\gamma_2) = (\gamma_2 J)_J$. Repeating with $\gamma_1,\gamma_2$ swapped yields $\widehat{\beta} \in \mathrm{End}(\widehat{P})$ mapping $i(\gamma_2)$ to $i(\gamma_1)$.
\end{description} 
\end{proof}

\end{document}
